%!TEX root = ../thesis.tex
%*******************************************************************************
%****************************** Third Chapter **********************************
%*******************************************************************************
\chapter{Preliminaries}

% **************************** Define Graphics Path **************************
\ifpdf
    \graphicspath{{Chapter2/Figs/Raster/}{Chapter2/Figs/PDF/}{Chapter2/Figs/}}
\else
    \graphicspath{{Chapter2/Figs/Vector/}{Chapter2/Figs/}}
\fi

\section{Basic Definitions}
In this section, I introduce the definitions and notations used in this work.
\newline

\subsubsection{Elections}

\noindent\textbf{Approval Elections}
\indent
Election is a pair \(E = (C, V)\) where 
\(C = \{c_1,...,c_m\}\) is a set of candidates and 
\(V = \{v_1,...,v_n\}\) is a list of voters. 
Each voter \(v \in V\) approves a subset of the candidates and disapproves the remaining candidates,
I will be referring to it as a \mbox{\textit{approval set}} and for each voter $v \in V$ denote it as $A(v)$, when
using $A$ just by itself I refer to all \mbox{\textit{approval set}} in the given election.

\subsubsection{Committee Selection Rules}

Voting rules are algorithms used to select winners (committee) of a given election.
Each rule has different properties and must be selected accordingly to the goal of the election.
\newline\indent For the detailed overview of multiwinner voting rules I point to \cite{faliszewskiSkowronMultiWinnerOverview}
\newline

\noindent\textbf{Hamming Distance Minimax Approval Voting (MAV)}
\indent
Selects a committee $W$ that minimizes the maximum Hamming distance between \(W\) and voters' approval ballots.

For an election $E$ and a resulting committee of size $k$ I denote winning committee as $W = MAV(E,k)$
\indent\cite{minimax-Brams-Kilgour-Sanver}

\subsubsection{Domain restrictions}

Many voting rules with interesting properties are NP-hard, however by introducing
certain restrictions to the structure of election it is possible to find a polynomial algorithm.

\noindent\textbf{Single Peaked}
\indent
SPOC  - \cite{DBLP:journals/jair/PetersL20}

\noindent\textbf{Single Crossing}
\indent

\noindent\textbf{VR / CR}
\indent
\cite{DBLP:journals/corr/ElkindL15}

\newpage
\section{(WIP) Related Work - TODO : Should Related Work section be completely removed?}

\begin{itemize}
    \item An Analysis of Approval-Based Committee Rules for 2D-Euclidean Elections by \cite{falisz}
    \item Overview of various domain restriction in Structure in Dichotomous Preferences by \cite{DBLP:journals/corr/ElkindL15}
    \item Minimax Approval Voting rule introduced by \cite{minimax-Brams-Kilgour-Sanver}
    \item Polynomial algorithms for MAV in VR and CR domains by \cite{mav}
    \item Example of a new domain restriction (SPOC) and recognition algorithms by \cite{DBLP:journals/jair/PetersL20}
    \item Example of a NP-hard domain recognition algorithms and theorems by \cite{peters2016recognising}
    \item Alternative approach to C1P for domain recognition by \cite{DBLP:journals/jair/MagieraF19}
\end{itemize}

